\subsection{Knowledge Distillation}
\begin{foldable}
  Transferring knowledge from a large model (teacher) to a smaller model (student) was first introduced in~\cite{bucila2006model}, then formally popularized in~\cite{hinton2015distilling}.
  Its key idea is to exploit the knowledge encoded in the teacher's `softened' probabilistic outputs, which are continuous and often more informative than the `hard' class labels.
  The additional supervision helps the student capture underlying inter-class relationships and improve its performance.
\end{foldable}

\begin{foldable}
  Existing KD methods often involve logits-based, feature-based, and relation-based knowledge~\cite{gou2021knowledge}.
  Logits-based methods use the teacher's logits (\ie, pre-softmax activations) as the guidance for the student~\cite{ba2014deep,hinton2015distilling,nguyen2021knowledge}, optionally with adaptive temperature~\cite{li2023curriculum} or an ensemble of teachers~\cite{guo2020online}.
  Feature-based methods extract knowledge from the teacher's internal attributes \eg, parameters and inter-layer connections~\cite{liu2019knowledge}, intermediate features~\cite{romero2014fitnets}, attention maps~\cite{zagoruyko2016paying}, and feature-space probability distribution~\cite{passalis2018learning}.
  Finally, relation-based methods distill the knowledge in the relationships between different layers or data samples~\cite{park2019relational,yim2017gift,chen2020learning}.
  Although useful, these methods would struggle when only a few images are available in the distillation set and we can only access the black-box teacher's final predictions.
\end{foldable}

\subsection{Knowledge distillation with limited data.}
\begin{foldable}
  The fundamental challenge of distillation on limited data is that the few available images cannot fully represent the actual images manifold.
  Consequently, the student fails to generalize to an unseen set.
  Thus, few-shot KD methods leverage data augmentation or generative models to augment the distillation set~\cite{kimura2018few,kong2020learning,wang2020neural,nguyen2022black}.
  Their approaches also vary depending on the accessibility of the teacher.
  FSKD~\cite{kimura2018few} uses trainable synthetic images that are part of the student's parameters.
  These images require adversarial signals from a white-box teacher for optimization.
  WaGe~\cite{kong2020learning} explores the data space in the approximate neighborhood of the few-shot data samples.
  This requires the logits of a white-box teacher for a Wasserstein-based loss.
  Most recently, BBKD~\cite{wang2020neural} and FS-BBT~\cite{nguyen2022black} only require a black-box teacher for the distillation.
  Both use MixUp to construct synthetic images for a larger distillation set for the student.
  While BBKD embraces active learning to choose synthetic images for distillation, FS-BBT trains a CVAE as an additional source of synthetic images.
  FS-BBT is the current SOTA few-shot KD method.
\end{foldable}

\subsection{Data-free knowledge distillation.}
\begin{foldable}
  As an even more extreme setting than few-shot, data-free KD methods assume no training data to train the student.
  To compensate, they often extract the knowledge from a white-box teacher \eg, its activation statistics~\cite{lopes2017data}, gradients~\cite{nayak2019zero,chen2019data}, and batch-norm layers~\cite{yin2020dreaming}.
  So far, only two methods use black-box teachers, namely ZSDB3KD~\cite{wang2021zero} and IDEAL~\cite{zhang2022ideal}.
  Although the data-free KD setting has several applications \eg, model inversion attack~\cite{zhang2020secret}, it is not really practical in the context of model compression.
  When we distill a large model, we often have data to validate the large and distilled models.
  Thus, few images can be also used for the distillation phase.
  As will be shown in \S\ref{ssec:04-experiments-comparison_zero_shot}., even with a small portion of data, our method significantly outperforms data-free KD methods.
\end{foldable}

\subsection{Generative Adversarial Networks in Knowledge Distillation}
\begin{foldable}
  Using Generative Adversarial Networks (GAN) to generate synthetic images is a popular approach in \textit{data-free} KD methods~\cite{chen2019data,addepalli2020degan,do2022momentum,zhang2022ideal}.
  However, our WGAN is trained with a novel scheme with \textit{three-fold differences}:
  \begin{enumerate}
    \item {
        \textbf{Generator:} Existing methods typically rely on the teacher's signals to tame the generator with extra losses \eg, one-hot, feature, and diversity losses.
        However, these losses are strictly inapplicable in our setting due to the absence of gradients from a black-box teacher.
      }
    \item {
        \textbf{Discriminator:} Existing methods cannot use a discriminator due to the lack of real images.
        In contrast, since our method allows to use any available real images, it uses the discriminator to distinguish real/fake images.
        Moreover, our method also trains the discriminator with high-confidence images to boost the diversity in image generation.
      }
    \item {
        \textbf{Teacher:} The teacher is trained for predicting image labels rather than distinguishing real/fake images.
        However, existing methods usually force the teacher to a discriminator's role, exposing it to knowledge gap.
        In contrast, we use teacher in a way where its intended strength is guaranteed to support our WGAN training.
      }
  \end{enumerate}
\end{foldable}
